\chapter{Communication véhicule-à-véhicule (V2V)}
\label{chap:v2v}
% ⭐ TON chapitre vedette : contribution conçue et réalisée de bout en bout
% par toi, from scratch. Rédige-le en « j'ai choisi / j'ai conçu / j'ai
% démontré ».

\section{Contexte}
% TODO : pour la sécurité coopérative, les véhicules doivent s'annoncer
% mutuellement (CAM) et s'alerter (DENM) — normes ETSI vues au chapitre 2.

\section{Problématique}
% TODO : implémenter une communication V2V conforme ETSI sur du matériel
% standard (WiFi, sans carte 802.11p dédiée), intégrée à ROS 2, et la
% démontrer entre deux véhicules réels.

\section{Mission 1 : Mise en \oe uvre de la pile Vanetza embarquée}
\objectif
\methodes
% TODO : choix de Vanetza (comparaison avec les alternatives), compilation
% croisée/embarquée avec backend OpenSSL, socktap sur interface WiFi.
\resultats
\discussion

\section{Mission 2 : Émission et réception de CAM intégrées à ROS 2}
\objectif
\methodes
% TODO : v2v_node — contenu des CAM, fréquence, topics /v2v/remote_vehicles,
% /v2v/alert ; stratégies de positionnement (GPS simulé, odométrie, distance
% lidar réelle via lidar_neighbor).
\resultats
% TODO : captures des CAM reçus, portée, fréquence effective.
\discussion

\section{Mission 3 : Alertes DENM déclenchées par la perception}
\objectif
\methodes
% TODO : chaînage perception (chap. 6) -> DENM : détection piéton YOLO
% (/obstacle/detections) -> émission DENM -> réaction du véhicule récepteur.
\resultats
% TODO : LA démo à deux véhicules : scénario, mesures, captures d'écran.
\discussion

\section{Conclusion}
